\section{Several founders}
\label{dist:sec:moi}

The single-cell law of \cref{dist:sec:geometric,dist:sec:burst-size} was
derived from one founder. The first probe of how far it reaches asks what
happens when there are several.

In applications a single founder is not always the right picture: a
macrophage may be invaded by several bacteria at once, and in the laboratory
the multiplicity of infection is a control parameter. Nothing essential
changes --- the process is still linear and $G_k=G^k$ still holds --- but the
bookkeeping does, and two plausible guesses turn out to be wrong. Throughout,
$k$ founders share one cell and hence one catastrophe clock: rupture occurs at
rate $\delta$ times the \emph{total} load, and dumps the whole cell.

\subsection{Fixation functions for \texorpdfstring{$k$}{k} founders}
\label{dist:sec:moi-fixation}

By the branching property the $k$ founder lineages evolve independently until
the shared clock rings. No burst occurs by $t$ precisely when no lineage has
burst; internal extinction occurs precisely when every lineage has cleared
internally. Hence
\begin{equation}
I_k(t)=I(t)^k,\qquad
D_k(t)=D(t)^k,\qquad
\Ifix[,{k}](t)=I(t)^k-D(t)^k.
\label{dist:eq:ihatk}
\end{equation}

\begin{remark}[A guess that fails]
\label{dist:rem:ihatk}
The tempting guess is $\Ifix[,{k}]=\Ifix^{\,k}$: each lineage stays productive
with probability $\Ifix$, independently. But ``productive'' for the whole
cell means \emph{not fixed}, and the intersection of the independent events
``lineage $i$ avoids bursting and does not go extinct'' is not the event
``the cell is not fixed'': the cell also stops being productive when every
lineage dies out. Correctly, $\Ifix[,{k}]=I^k-D^k$, whereas
$\Ifix^k=(I-D)^k$; already for $k=2$ the two differ by
$I^2-D^2-(I-D)^2=2D(I-D)\neq0$ whenever $\mu>0$. Numerically, at
$(\lambda,\mu,\delta)=(1,0.2,0.05)$ and $t=1$: the correct
$\Ifix[,{2}]=I^2-D^2\approx0.8458$, while $\Ifix^2\approx0.6542$. When $\mu=0$,
$D\equiv0$ and the two coincide --- again the invisible case.
\end{remark}

\subsection{Moments and the release kernel}
\label{dist:sec:moi-moments}

The moments fall out of the $k$-founder PGF of
\cref{dist:sec:pgf-k-founders}, $G_k=G^k$, where $G$ is the single-founder
defective PGF. Differentiating at $z=1$, with $G(1,t)=I(t)$,
$\partial_zG|_{z=1}=J(t)$ and $\partial_z^2G|_{z=1}=K(t)-J(t)$,
\begin{equation}
J_k(t)=k\,I(t)^{k-1}J(t),
\label{dist:eq:jk}
\end{equation}
and, from the second factorial moment,
\begin{align}
K_k(t)&=\partial_z^2G_k\big|_{z=1}+\partial_zG_k\big|_{z=1}\nonumber\\
&=k(k-1)I(t)^{k-2}J(t)^2+kI(t)^{k-1}\lrb{K(t)-J(t)}+kI(t)^{k-1}J(t)\nonumber\\
&=k\,K(t)I(t)^{k-1}+k(k-1)J(t)^2I(t)^{k-2}.
\label{dist:eq:kk}
\end{align}
Multiplying by the catastrophe rate turns $K_k$ into the expected release
flux --- the second item of the consignment of \cref{dist:sec:exports}, and
the kernel on which the renewal model of \ChPop{} is built:
\begin{equation}
g_k(t)=\delta K_k(t)
=\delta\lrb{k\,K(t)I(t)^{k-1}+k(k-1)J(t)^2I(t)^{k-2}}.
\label{dist:eq:gkflux}
\end{equation}
\Cref{dist:fig:moi}(a) plots $g_k$ against age for $k=1,\ldots,4$. The
ordering in $k$ is superlinear early and reverses late, so the enhancement is
a property of the cell's age rather than a fixed multiplier.

\begin{remark}[The free-sum trap]
\label{dist:rem:gk}
It is tempting to write $g_k=\delta\lrb{kK+k(k-1)J^2}$, the second moment of
a free sum of $k$ independent copies. That form omits the powers of $I$: the
second-moment contribution of a configuration is counted only while the
shared clock has \emph{not} rung. At $(\lambda,\mu,\delta)=(1,0.2,0.05)$,
$t=1$, $k=2$: the free-sum form gives $K_2\approx23.39$, while
\cref{dist:eq:kk} gives $K_2\approx22.264$, which agrees with a direct
master-equation computation. The difference is the superlinearity of $g_k$
in $k$ wearing a disguise: even the free-sum form is superlinear, but the
true kernel is superlinear \emph{and} tempered by survival. That a cell's
expected release flux depends superlinearly on its founding multiplicity is
a genuine prediction of the model --- one that classical constant-$p$ models
cannot express, and that low-multiplicity single-cell assays can test.
\end{remark}

\subsection{Lifetime yield}
\label{dist:sec:moi-yield}

The lifetime mean release from $k$ founders is
$V_\infty^{(k)}=\int_0^\infty g_k(t)\,\dd t$. Carrying out the integration,
which is deferred to \cref{dist:app:vk},
\begin{equation}
V_\infty^{(1)}=V_\infty=\frac{a(1-b)}{a-1},
\label{dist:eq:vk1}
\end{equation}
and, for $k\ge2$,
\begin{equation}
V_\infty^{(k)}
=(1-b^k)
+\frac{2\lambda}{\delta}\Biggl[\frac{1}{k+1}-b^k+\frac{k\,b^{k+1}}{k+1}\Biggr]
+k(k-1)\frac{\lambda}{\delta}\int_1^b (x-a)(x-b)\,x^{k-2}\,\dd x.
\label{dist:eq:vk}
\end{equation}
In the no-death case the answer collapses:
\begin{equation}
V_\infty^{(k)}=k+\frac{\lambda}{\delta}\qquad(\mu=0).
\label{dist:eq:vkmuzero}
\end{equation}
Indeed, when $\mu=0$ the event type is independent of the load
(\cref{dist:sec:key-observation}), so the number of births before the single
catastrophe is $\mathrm{Geom}_0(s)$ with mean $\rho/s=\lambda/\delta$
whatever the founding load, and the total release is simply $k$ founders plus
that many births. With $\lambda=1$ and $\delta=0.1$, for instance,
$V_\infty^{(k)}=11,12,13,14,15,\ldots$

Two consequences are worth pausing on, and they pull in opposite directions.
The yield is \emph{not} $kV_\infty$: the first catastrophe dumps the entire
cell once, so additional founders add their own expected descendants but not
another cell's worth of release. At $(\lambda,\mu,\delta)=(1,0.2,0.05)$ a
naive scaling would predict $2V_\infty\approx27.97$ for two founders; the
true value is $V_\infty^{(2)}\approx17.43$, and
\cref{dist:fig:subextensive} shows the gap widening in $k$. Yet the
\emph{conditional} mean burst size given eventual burst,
$V_\infty^{(k)}/(1-b^k)$, \emph{increases} with $k$ --- $17.23$, $18.07$,
$19.02$, $21.00$ for $k=1,2,3,5$ at the same parameters. Multiply infected
cells are more likely to burst at all, and when they do, they burst bigger.
The two statements are consistent because the first is unconditional and the
second is not, which is \cref{dist:tab:threemeans}'s distinction appearing
again in a new place.

The geometric structure of \cref{dist:thm:geometric} survives all of this
untouched, the load in a multiply infected cell being a sum of independent
copies of the single-founder load. Multiplicity changes the bookkeeping, not
the shape; \cref{dist:sec:chained} finds the boundary the other way.

\begin{figure}[tbp]
\centering
\includegraphics[width=0.68\textwidth]{fig12_yield_subextensive/fig12}
\caption{Unconditional mean yield $V_\infty^{(k)}$ against founding
multiplicity $k$, compared with the naive independent-cell line $kV_\infty$,
for $(\lambda,\mu,\delta)=(1,0.2,0.05)$. The true yield is sub-extensive
because all founders share one cell and hence one catastrophe; for $k=2$,
$V_\infty^{(2)}=17.43$ rather than $2V_\infty=27.97$.}
\label{dist:fig:subextensive}
\end{figure}

\begin{figure}[tbp]
\centering
\includegraphics[width=0.9\textwidth]{fig13_moi/fig13}
\caption{Multiplicity of infection at $(\lambda,\mu,\delta)=(1,0.2,0.05)$.
(a) Release fluxes $g_k(t)=\delta K_k(t)$ from \cref{dist:eq:gkflux} for
$k=1,\ldots,4$: early ages are ordered superlinearly in founding load, while
late fluxes reverse as multi-founder lineages fix faster, so the enhancement
is age-dependent rather than a pointwise ordering valid for all $t$.
(b) Conditional mean burst $V_\infty^{(k)}/(1-b^k)$ increases with $k$; the
dashed line is the $k=1$ value $a/(a-1)=17.23$.}
\label{dist:fig:moi}
\end{figure}
